\section{Introduction}
A generative model of a random field should describe one probabilistic object across the resolutions at which that object is observed. A simulator of an order book, for example, may produce liquidity aggregated over price intervals, liquidity at individual prices, and latent reserve quantities. A model of a physical field may produce low Fourier modes and their unresolved fluctuations. In both cases, fitting a collection of finite arrays leaves open whether those arrays belong to a common law and whether the dynamics used to generate them respect the chosen observation hierarchy.

Neural operators provide architectures for approximating maps between function spaces \citep{kovachki2023}. Their generative counterparts learn a pushforward law rather than a single solution operator. Generative adversarial neural operators make this construction explicit \citep{rahman2022}. Diffusion and flow formulations extend the available transports to function-valued random objects \citep{pidstrigach2025,lim2025,kerrigan2024functional}. These developments establish the basic modeling language. The present paper studies a compatibility problem within that language: how should a generative transport respond to the removal or addition of observable coordinates?

Four notions must be distinguished. \emph{Discretization convergence} asks whether increasingly accurate numerical realizations approach a continuum operator. \emph{Projective consistency} asks whether the projected fine law equals the coarse law. \emph{Pathwise consistency} strengthens this to equality of projected samples under a specified coupling, such as shared coarse randomness. \emph{Autonomous projected dynamics} asks whether the velocity or transition law of the coarse state is determined by that state alone. Autonomous projected dynamics can fail even when a continuum model exists and every coarse marginal is known exactly.

Our starting point is an information identity. For a full velocity $v_s(X_s)$ and an orthogonal observation $P$, the best coarse-only approximation to $Pv_s(X_s)$ is its conditional expectation given $PX_s$. The remaining conditional variance is an irreducible population error. We apply this identity to a prescribed generative bridge, evaluate it exactly for correlated Gaussian targets, and show why projecting the optimal velocity can change the generated endpoint law. The obstruction concerns the specified bridge and the chosen information restriction; it does not rule out other transports with the same target law.

The constructive result uses conditional refinement. After generating a coarse sample, the model holds it fixed and samples the next orthogonal detail band conditional on it. This is a resolution hierarchy of conditional transports. If $\eps_j$ is the conditional distribution error at level $j$, $L_j$ is the sensitivity of the learned detail law to its input context, and $d_j$ is a propagated error bound, then
\begin{equation}\label{eq:intro-recurrence}
 d_j^2=d_{j-1}^2+(\eps_j+L_jd_{j-1})^2,
 \qquad d_{-1}=0.
\end{equation}
The recurrence is sharp in the kernel class considered here. We also compute its exact worst-case amplification of aggregate fitting energy through a two-dimensional eigenvalue recursion. This gives a necessary and sufficient squared-sensitivity condition for uniform robustness over the full conditional-kernel class. Orthogonality also yields the exact decomposition
\begin{equation}\label{eq:intro-tail}
 W_2^2(\mu,\widehat\mu_m)
 =\E\norm{(I-P_m)U}^2+W_2^2((P_m)\push\mu,\widehat\mu_m)
\end{equation}
when $\widehat\mu_m$ is supported on the retained subspace. Consequently, conditional fitting errors and unresolved tail energy enter through different mathematical mechanisms.

The results extend beyond a fixed truncation. If $\sum_j\eps_j^2<\infty$ and $\sum_jL_j^2<\infty$, the refinement construction defines a genuine Hilbert-valued random element, with an explicit Wasserstein bound independent of the terminal level. Conditional flow-matching excess losses provide the $\eps_j$ through a one-sided ODE stability estimate. This connects a population training objective to both the existence and the accuracy of the function-space law.

The same framework separates geometry, architecture, and statistical information. Stable Riesz coordinates or a bi-Lipschitz decoder transfer coefficient certificates to nonorthogonal physical representations. A coordinatewise dependency bound can improve on the class-uniform recurrence when later kernels ignore inaccurate earlier coordinates. Finally, an observed-sample audit controls scalar conditional laws through empirical transport in feature cells and an independent aggregation split. It removes the need to evaluate a target conditional mean, while exposing the sufficient-context and regularity assumptions that make conditional estimation possible. An unseen-cell lower bound explains why these assumptions cannot be omitted.

\subsection{Relation to established transport and operator theory}
The construction of function-space neural operators and their approximation properties is developed by \citet{kovachki2023}. Continuum attention provides a corresponding operator interpretation of transformer layers \citep{calvello2024}. Flow matching and stochastic interpolants identify velocity fields by regression on prescribed probability paths \citep{lipman2023,albergo2023}. Functional flow matching extends that viewpoint to infinite-dimensional laws \citep{kerrigan2024functional}. Valid reference measures and endpoint regularity remain substantive assumptions; trace-class noise alone does not settle the regularity of the sampling equation. \citet{chen2026} study how spectral noise and interpolation schedules affect these issues.

Conditional transport is also established. \citet{hosseini2025} develop triangular optimal transport on separable function spaces. \citet{kerrigan2024conditional} give a dynamical formulation and simulation-free conditional flows. \citet{chemseddine2025} study conditional Wasserstein distances and their relevance for learning conditional laws. Our conditional kernel discrepancy is an instance of this geometry. Triangular Gaussian cost comparisons and adapted Gaussian transport are treated by \citet{bukin2018,acciaio2025}. We use standard disintegration and kernel randomization \citep{kallenberg2021}; we do not claim a new existence theorem for conditional transport itself.

The contribution is the connection between a prescribed-bridge autonomy defect, a sharp orthogonal error recurrence for an entire resolution hierarchy, and a summability criterion that converts conditional fitting guarantees into a coherent function-space law. The finite-dimensional Gaussian commutation criterion makes the obstruction explicit. A companion Gaussian result computes the exact transport-cost premium of resolution autonomy and constructs an autonomous bridge attaining it. The observation minimax result and physical-time kernel criterion delimit what the resulting law can identify and what an autonomous simulator can preserve. Each statement is proved under explicit assumptions. The numerical section tests these statements and trains a shared spectral neural generator on noisy field coefficients. An architecture theorem provides its global sensitivity and energy bounds, and an independent conditional audit supplies simultaneous finite-sample law certificates. The controlled experiment establishes performance within its stated synthetic model.

The closest perturbation antecedents include \citet{rudolf2018}, who control the effect of one-step kernel errors under Wasserstein stability and ergodicity assumptions. Our hierarchy is not an ergodic chain on a fixed state space: it retains the entire previous state and appends an orthogonal coordinate. The resulting improvement is displayed in \Cref{sec:comparison}. \citet{irons2022} establish statistical consistency and rates for triangular-flow estimators under smoothness assumptions. Their estimator theory and our conditional validation certificate address different tasks; the latter certifies a fixed trained hierarchy using an independent audit, without asserting an optimization rate.

\subsection{Organization}
\Cref{sec:setting} fixes the probabilistic and geometric setting. \Cref{sec:obstruction} derives the autonomy defect. \Cref{sec:refinement,sec:certificate} construct conditional refinement and prove the main error and existence results. \Cref{sec:fm} links the certificate to flow matching. \Cref{sec:information} treats observation limits and physical time. \Cref{sec:riesz} treats nonorthogonal and nonlinear representations. \Cref{sec:architecture,sec:observed-audit} give implemented stability, observed-data auditing, and dependency-sensitive propagation. \Cref{sec:experiments,sec:application} report numerical evidence and a market-modeling application.
